\documentclass[10pt]{extarticle}
\usepackage[a4paper,margin=1.5cm]{geometry}
\usepackage[T1]{fontenc}
\usepackage{amsmath,amssymb}
\usepackage{graphicx}
\usepackage[export]{adjustbox}
\usepackage{multicol}
\usepackage{enumitem}
\usepackage{booktabs,array,tabularx}
\setlength{\columnsep}{0.7cm}
\setlength{\parindent}{0pt}
\setlength{\parskip}{3pt}
\pagestyle{empty}
\makeatletter
% figure registry (figdefs.tex is generated by the extract step)
\newcommand{\qfigdef}[4]{\@namedef{qf@file@#1}{#2}\@namedef{qf@w@#1}{#3}\@namedef{qf@wide@#1}{#4}}
\newcommand{\qf@check}[1]{\@ifundefined{qf@file@#1}{\PackageError{score2}{Unknown figure id #1}{}}{}}
\newcommand{\qf@img}[4][]{\includegraphics[width=#2\dimexpr\@nameuse{qf@w@#4}pt\relax,max width=#3,max height=0.45\textheight,#1]{\@nameuse{qf@file@#4}}}
% Full-width questions: finalize wraps every question that contains a wide figure in \qfullbegin...\qfullend,
% so the WHOLE question (stem, figures, options) is set across the page in reading order, then the two columns
% resume. The question's head - its number, stem and first figure (\qfullsplit marks the end of the head) -
% is kept on one page, so a number is never left at a page foot with its figure overleaf; the rest may break
% across pages normally, so no big gaps are left. Without a \qfullsplit, the whole question is kept together
% if it is under 45% of a page. Never write these yourself.
\newif\ifqf@full
\newif\ifqf@boxed
\newsavebox{\qf@box}
\newcommand{\qf@open}{\global\qf@fulltrue\par\global\qf@boxedtrue
  \setbox\qf@box\vbox\bgroup\hsize\textwidth\linewidth\textwidth}
\newcommand{\qf@emit}[1]{\par\egroup\global\qf@boxedfalse
  \ifdim\dimexpr\ht\qf@box+\dp\qf@box\relax<#1\textheight\noindent\usebox{\qf@box}\par\else\unvbox\qf@box\fi}
\newcommand{\qfullsplit}{\ifqf@boxed\qf@emit{0.85}\fi}
\newcommand{\qf@close}{\ifqf@boxed\qf@emit{0.45}\fi\global\qf@fullfalse}
\newcommand{\qfullbegin}{\end{multicols}\qf@open}
\newcommand{\qfullend}{\qf@close\begin{multicols}{2}}
\newcommand{\qfullstart}{\qf@open}
\newcommand{\qfulldone}{\qf@close}
% \qfig[scale]{id}: block figure, centred, directly where it is written (size set by the tool).
% Fallback if a wide figure is not inside a full-width question: break out of the columns in place - no floats.
\newcommand{\qfig}[2][1]{\qf@check{#2}%
  \ifqf@full\qf@blk{#1}{#2}\else
  \ifnum\@nameuse{qf@wide@#2}=1 \expandafter\@firstoftwo\else\expandafter\@secondoftwo\fi
  {\end{multicols}\begin{center}\qf@img{#1}{\textwidth}{#2}\end{center}\begin{multicols}{2}}%
  {\qf@blk{#1}{#2}}\fi}
\newcommand{\qf@blk}[2]{\par\smallskip{\centering\qf@img{#1}{\linewidth}{#2}\par}\smallskip}
% \qfigi[scale]{id}: inline figure (inside option lists / tables / minipages); never breaks columns.
\newcommand{\qfigi}[2][1]{\qf@check{#2}\qf@img[valign=c]{#1}{\linewidth}{#2}}
\makeatother
% structure macros - lint reads these to check completeness and figure placement
\newcommand{\papersec}[1]{\par\bigskip{\centering\Large\bfseries PAPER-#1\par}\medskip}
\newcommand{\partsec}[2]{\par\medskip\noindent{\large\bfseries PART-#1: #2}\par\smallskip}
\newcommand{\sect}[1]{\par\smallskip\noindent\textbf{#1}\par}
\newcommand{\qn}[1]{\par\smallskip\noindent\textbf{#1.}~}
\newcommand{\paraq}[2]{\par\smallskip\noindent\textbf{Paragraph for Questions #1 and #2}\par}
\newenvironment{qoptions}{\begin{itemize}[leftmargin=1.8em,itemsep=1pt,topsep=2pt,label={}]}{\end{itemize}}
\newcommand{\opt}[2]{\item[(#1)] #2}
% figure registry (generated by score2_tools.py finalize)
\qfigdef{p1_phy_s1_q01_1}{TOAS - Test 18 (30 Oct 2022)-img/p1_phy_s1_q01_1.png}{119.0}{0}
\qfigdef{p1_phy_s1_q03_1}{TOAS - Test 18 (30 Oct 2022)-img/p1_phy_s1_q03_1.png}{129.4}{0}
\qfigdef{p1_phy_s1_q04_1}{TOAS - Test 18 (30 Oct 2022)-img/p1_phy_s1_q04_1.png}{142.8}{0}
\qfigdef{p1_phy_s1_q04_A}{TOAS - Test 18 (30 Oct 2022)-img/p1_phy_s1_q04_A.png}{93.5}{0}
\qfigdef{p1_phy_s1_q04_B}{TOAS - Test 18 (30 Oct 2022)-img/p1_phy_s1_q04_B.png}{93.1}{0}
\qfigdef{p1_phy_s1_q04_C}{TOAS - Test 18 (30 Oct 2022)-img/p1_phy_s1_q04_C.png}{93.5}{0}
\qfigdef{p1_phy_s1_q04_D}{TOAS - Test 18 (30 Oct 2022)-img/p1_phy_s1_q04_D.png}{88.4}{0}
\qfigdef{p1_phy_s1_q05_1}{TOAS - Test 18 (30 Oct 2022)-img/p1_phy_s1_q05_1.png}{87.2}{0}
\qfigdef{p1_phy_s2_q08_1}{TOAS - Test 18 (30 Oct 2022)-img/p1_phy_s2_q08_1.png}{92.5}{0}
\qfigdef{p1_phy_s2_q08_2}{TOAS - Test 18 (30 Oct 2022)-img/p1_phy_s2_q08_2.png}{110.6}{0}
\qfigdef{p1_phy_s2_q08_3}{TOAS - Test 18 (30 Oct 2022)-img/p1_phy_s2_q08_3.png}{79.4}{0}
\qfigdef{p1_phy_s2_q08_4}{TOAS - Test 18 (30 Oct 2022)-img/p1_phy_s2_q08_4.png}{87.2}{0}
\qfigdef{p1_phy_s2_q09_1}{TOAS - Test 18 (30 Oct 2022)-img/p1_phy_s2_q09_1.png}{83.5}{0}
\qfigdef{p1_phy_s2_q09_2}{TOAS - Test 18 (30 Oct 2022)-img/p1_phy_s2_q09_2.png}{83.5}{0}
\qfigdef{p1_phy_s2_q09_3}{TOAS - Test 18 (30 Oct 2022)-img/p1_phy_s2_q09_3.png}{85.2}{0}
\qfigdef{p1_phy_s2_q09_4}{TOAS - Test 18 (30 Oct 2022)-img/p1_phy_s2_q09_4.png}{88.2}{0}
\qfigdef{p1_phy_s2_q10_1}{TOAS - Test 18 (30 Oct 2022)-img/p1_phy_s2_q10_1.png}{92.3}{0}
\qfigdef{p1_phy_s2_q10_2}{TOAS - Test 18 (30 Oct 2022)-img/p1_phy_s2_q10_2.png}{87.6}{0}
\qfigdef{p1_phy_s2_q10_3}{TOAS - Test 18 (30 Oct 2022)-img/p1_phy_s2_q10_3.png}{86.9}{0}
\qfigdef{p1_phy_s2_q10_4}{TOAS - Test 18 (30 Oct 2022)-img/p1_phy_s2_q10_4.png}{97.5}{0}
\qfigdef{p1_phy_s3_q03_1}{TOAS - Test 18 (30 Oct 2022)-img/p1_phy_s3_q03_1.png}{109.4}{0}
\qfigdef{p1_phy_s3_q05_1}{TOAS - Test 18 (30 Oct 2022)-img/p1_phy_s3_q05_1.png}{161.1}{0}
\qfigdef{p1_phy_s3_q06_1}{TOAS - Test 18 (30 Oct 2022)-img/p1_phy_s3_q06_1.png}{129.3}{0}
\qfigdef{p1_chem_s1_q04_1}{TOAS - Test 18 (30 Oct 2022)-img/p1_chem_s1_q04_1.png}{98.9}{0}
\qfigdef{p1_chem_s1_q05_1}{TOAS - Test 18 (30 Oct 2022)-img/p1_chem_s1_q05_1.png}{153.0}{0}
\qfigdef{p1_chem_s1_q06_1}{TOAS - Test 18 (30 Oct 2022)-img/p1_chem_s1_q06_1.png}{230.6}{0}
\qfigdef{p1_chem_s2_q09_1}{TOAS - Test 18 (30 Oct 2022)-img/p1_chem_s2_q09_1.png}{122.4}{0}
\qfigdef{p1_chem_s2_q09_5}{TOAS - Test 18 (30 Oct 2022)-img/p1_chem_s2_q09_5.png}{33.8}{0}
\qfigdef{p1_chem_s2_q09_2}{TOAS - Test 18 (30 Oct 2022)-img/p1_chem_s2_q09_2.png}{124.2}{0}
\qfigdef{p1_chem_s2_q09_6}{TOAS - Test 18 (30 Oct 2022)-img/p1_chem_s2_q09_6.png}{27.9}{0}
\qfigdef{p1_chem_s2_q09_3}{TOAS - Test 18 (30 Oct 2022)-img/p1_chem_s2_q09_3.png}{108.0}{0}
\qfigdef{p1_chem_s2_q09_7}{TOAS - Test 18 (30 Oct 2022)-img/p1_chem_s2_q09_7.png}{38.0}{0}
\qfigdef{p1_chem_s2_q09_4}{TOAS - Test 18 (30 Oct 2022)-img/p1_chem_s2_q09_4.png}{144.5}{0}
\qfigdef{p1_chem_s2_q09_8}{TOAS - Test 18 (30 Oct 2022)-img/p1_chem_s2_q09_8.png}{26.5}{0}
\qfigdef{p1_chem_s2_q09_9}{TOAS - Test 18 (30 Oct 2022)-img/p1_chem_s2_q09_9.png}{29.2}{0}
\qfigdef{p1_chem_s3_q06_1}{TOAS - Test 18 (30 Oct 2022)-img/p1_chem_s3_q06_1.png}{150.3}{0}
\qfigdef{p1_chem_s3_q07_1}{TOAS - Test 18 (30 Oct 2022)-img/p1_chem_s3_q07_1.png}{242.4}{0}
\qfigdef{p1_chem_s3_q07_2}{TOAS - Test 18 (30 Oct 2022)-img/p1_chem_s3_q07_2.png}{221.3}{0}
\qfigdef{p1_chem_s3_q08_1}{TOAS - Test 18 (30 Oct 2022)-img/p1_chem_s3_q08_1.png}{203.8}{0}
\qfigdef{p1_chem_s3_q08_2}{TOAS - Test 18 (30 Oct 2022)-img/p1_chem_s3_q08_2.png}{214.3}{0}
\qfigdef{p1_chem_s3_q08_4}{TOAS - Test 18 (30 Oct 2022)-img/p1_chem_s3_q08_4.png}{151.6}{0}
\qfigdef{p2_phy_s1_q02_A}{TOAS - Test 18 (30 Oct 2022)-img/p2_phy_s1_q02_A.png}{75.7}{0}
\qfigdef{p2_phy_s1_q02_B}{TOAS - Test 18 (30 Oct 2022)-img/p2_phy_s1_q02_B.png}{75.2}{0}
\qfigdef{p2_phy_s1_q02_C}{TOAS - Test 18 (30 Oct 2022)-img/p2_phy_s1_q02_C.png}{70.5}{0}
\qfigdef{p2_phy_s1_q02_D}{TOAS - Test 18 (30 Oct 2022)-img/p2_phy_s1_q02_D.png}{75.2}{0}
\qfigdef{p2_phy_s1_q04_1}{TOAS - Test 18 (30 Oct 2022)-img/p2_phy_s1_q04_1.png}{82.0}{0}
\qfigdef{p2_phy_s2_q06_1}{TOAS - Test 18 (30 Oct 2022)-img/p2_phy_s2_q06_1.png}{93.7}{0}
\qfigdef{p2_phy_s2_q10_1}{TOAS - Test 18 (30 Oct 2022)-img/p2_phy_s2_q10_1.png}{108.3}{0}
\qfigdef{p2_phy_s3_q01_1}{TOAS - Test 18 (30 Oct 2022)-img/p2_phy_s3_q01_1.png}{159.9}{0}
\qfigdef{p2_phy_s3_q02_1}{TOAS - Test 18 (30 Oct 2022)-img/p2_phy_s3_q02_1.png}{119.5}{0}
\qfigdef{p2_phy_s3_q04_1}{TOAS - Test 18 (30 Oct 2022)-img/p2_phy_s3_q04_1.png}{119.6}{0}
\qfigdef{p2_phy_s3_q05_1}{TOAS - Test 18 (30 Oct 2022)-img/p2_phy_s3_q05_1.png}{146.5}{0}
\qfigdef{p2_phy_s3_q06_1}{TOAS - Test 18 (30 Oct 2022)-img/p2_phy_s3_q06_1.png}{105.3}{0}
\qfigdef{p2_phy_s3_q07_1}{TOAS - Test 18 (30 Oct 2022)-img/p2_phy_s3_q07_1.png}{112.7}{0}
\qfigdef{p2_chem_s1_q04_1}{TOAS - Test 18 (30 Oct 2022)-img/p2_chem_s1_q04_1.png}{224.4}{0}
\qfigdef{p2_chem_s2_q07_1}{TOAS - Test 18 (30 Oct 2022)-img/p2_chem_s2_q07_1.png}{161.1}{0}
\qfigdef{p2_chem_s2_q07_A}{TOAS - Test 18 (30 Oct 2022)-img/p2_chem_s2_q07_A.png}{81.8}{0}
\qfigdef{p2_chem_s2_q09_A}{TOAS - Test 18 (30 Oct 2022)-img/p2_chem_s2_q09_A.png}{130.1}{0}
\qfigdef{p2_chem_s2_q09_B}{TOAS - Test 18 (30 Oct 2022)-img/p2_chem_s2_q09_B.png}{161.5}{0}
\qfigdef{p2_chem_s2_q09_C}{TOAS - Test 18 (30 Oct 2022)-img/p2_chem_s2_q09_C.png}{173.7}{0}
\qfigdef{p2_chem_s2_q09_D}{TOAS - Test 18 (30 Oct 2022)-img/p2_chem_s2_q09_D.png}{134.7}{0}
\qfigdef{p2_chem_s2_q10_A}{TOAS - Test 18 (30 Oct 2022)-img/p2_chem_s2_q10_A.png}{79.3}{0}
\qfigdef{p2_chem_s2_q10_B}{TOAS - Test 18 (30 Oct 2022)-img/p2_chem_s2_q10_B.png}{138.9}{0}
\qfigdef{p2_chem_s2_q10_C}{TOAS - Test 18 (30 Oct 2022)-img/p2_chem_s2_q10_C.png}{123.9}{0}
\qfigdef{p2_chem_s2_q10_D}{TOAS - Test 18 (30 Oct 2022)-img/p2_chem_s2_q10_D.png}{92.5}{0}
\qfigdef{p2_chem_s3_q08_1}{TOAS - Test 18 (30 Oct 2022)-img/p2_chem_s3_q08_1.png}{169.0}{0}
\qfigdef{p2_chem_s3_q08_2}{TOAS - Test 18 (30 Oct 2022)-img/p2_chem_s3_q08_2.png}{137.4}{0}
\raggedcolumns
\makeatletter\@beginparpenalty=10000\makeatother
\newcolumntype{L}[1]{>{\raggedright\arraybackslash}p{#1}}
\begin{document}
\begin{center}{\Large\bfseries TOAS --- Test 18 (30 Oct 2022)}\end{center}
\begin{multicols}{2}
\raggedcolumns
\papersec{1}
\partsec{1}{PHYSICS}
\sect{SECTION-I(i)}
\qn{1} In given experiment on the photoelectric effect monochromatic plane wave of light falls on a metal plate of area 66.3 cm$^2$ at angle $30^\circ$ as shown. Electric field in light wave varies according to $E=E_0\cos(\omega t+\phi_0)$ where amplitude $E_0=15$ V/m, angular frequency is $\omega=9.5\times10^{15}$ rad/s. Dependence of photoelectric current $I$ on voltage applied $V$ between the electrodes is shown by graph. Select the correct statement (s).
\qfig{p1_phy_s1_q01_1}
\begin{qoptions}
\opt{A}{Work function of metal plate is 4.26 eV.}
\opt{B}{Power received be metal plate by radiation is $9.9\times10^{-4}$ W.}
\opt{C}{Quantum yield (photoelectric efficiency) in the experiment is 0.038.}
\opt{D}{If amplitude of electric field in incident light is made 25 V/m then saturation current will become 16.67 $\mu$A.}
\end{qoptions}
\qn{2} Select correct statement(s) :-
\begin{qoptions}
\opt{A}{If two particle having same mass have same De-Broglie wavelength then they have same kinetic energy and same momentum.}
\opt{B}{When a bomb at rest explodes into two unequal pieces then both pieces have same De-Brogli wavelength.}
\opt{C}{When a photon of wavelength $\lambda$ incident on a free electron at rest then it's De-Broglie wavelength becomes $\sqrt{\dfrac{h\lambda}{2mC}}$, where $m$ is mass of electron.}
\opt{D}{A proton and an alpha particle are moving along circular path under same uniform magnetic field, then ratio of radius will be 2 : 1 if they have same De-Broglie wavelength}
\end{qoptions}
\qn{3} A rod of length $r=3$ cm (exactly known) is pivoted about it's end and free to rotate in horizontal plane. A string of length $\ell=(20\pm1)$ cm is connected to opposite end of rod. A stone of mass $m$ is connected to other end of string. When rod rotate at rate $\omega=(7\pm0.07)$ rad/s then string makes angle $\theta=\left(37\pm\dfrac3\pi\right)^\circ$ with vertical. Then :
\qfig{p1_phy_s1_q03_1}
\begin{qoptions}
\opt{A}{mean value of acceleration due to gravity is 9.8 m/s$^2$}
\opt{B}{mean value of acceleration due to gravity is 12.25 m/s$^2$}
\opt{C}{\% error in measurement of $g$ is nearly 7.67\%}
\opt{D}{\% error in measurement of $g$ is nearly 5\%}
\end{qoptions}
\qn{4} A rigid wire loop of equilateral triangle shape having resistance $R$ is moving with constant velocity $v_0$ in plane of paper along $x$-axis. At $t=0$ edge $P$ of loop enter a region where there is a uniform magnetic field $B$ in to plane of paper. For sufficiently large is, the loop eventually enter the region. Let $x$ be the location of edge $P$. Let $\phi(x)$, $v(x)$, $I(x)$ and $F(x)$ represent magnetic flux through loop, velocity of loop, current in loop and force on the loop respectively as a function of $x$ counter-clockwise direction is taken as positive.

Which of following schematic plot(s) is (are) correct ? $\left[\dfrac{B^2L^3}{mR}=v_0\right]$
\qfig{p1_phy_s1_q04_1}
\begin{qoptions}
\opt{A}{\qfigi{p1_phy_s1_q04_A}}
\opt{B}{\qfigi{p1_phy_s1_q04_B}}
\opt{C}{\qfigi{p1_phy_s1_q04_C}}
\opt{D}{\qfigi{p1_phy_s1_q04_D}}
\end{qoptions}
\qn{5} A capacitor consists of two stationary plates shaped as a semi-circle of radius $R$ and a movable plate made of dielectric with permittivity $\in$ and capable of rotating about an axis $O$ between the stationary plates. The thickness of movable plate is $d$, which is nearly same as separation between stationary plates. Select correct statement(s). $\phi$ is angle of sector outside of capacitor.
\qfig{p1_phy_s1_q05_1}
\begin{qoptions}
\opt{A}{Capacitance of system is proportional to $\left\{\in-\dfrac\phi\pi(\in-\in_0)\right\}$}
\opt{B}{If potential difference $V$ is applied to capacitor electric torque about axis $O$ is $\dfrac{V^2R^2}{4d}(\in-\in_0)$ acting on movable plate.}
\opt{C}{If this capacitor is charged by charge $Q$ and then source is removed then electric torque about axis $O$ is $\dfrac{Q^2}{4(\in-\in_0)\,d}$.}
\opt{D}{If potential difference $V$ is applied to capacitor and system is released than movable plate will perform simple harmonic motion.}
\end{qoptions}
\qn{6} A thin uniform disk of mass $m$ and radius $R$ suspended by an elastic thread in horizontal plane in a liquid. The torsional constant of thread is $C$ and the resistance force acting per unit area of disk is $F=\eta v$ where $\eta$ is a constant and $v$ is the velocity of given element of disk. Take : $m=1$ kg, $R=0.5$ m, $C=12.5$ Nm/rad.
\begin{qoptions}
\opt{A}{Natural angular frequency of torsional oscillation is 10 rad/s}
\opt{B}{For $\eta\geqslant\dfrac{40}{\pi}\dfrac{\text{Ns}}{\text{m}^3}$, system will not be oscillatory.}
\opt{C}{For $\eta=\dfrac{24}{\pi}\dfrac{\text{Ns}}{\text{m}^3}$, system will perform damped oscillations with angular frequency 6 rad/s}
\opt{D}{For $\eta=\dfrac{24}{\pi}\dfrac{\text{Ns}}{\text{m}^3}$, Q-factor of system is $\dfrac56$.}
\end{qoptions}
\sect{SECTION-I(ii)}
\qn{7} Match list-I of nuclear process with list-II containing parent nucleus and one of end product of each process and then select the correct answer using codes given below the lists.
\begin{center}\footnotesize\setlength{\tabcolsep}{2pt}
\begin{tabular}{|L{0.8cm}L{1.6cm}|L{0.6cm}L{3.3cm}|}\hline
\multicolumn{2}{|c|}{List-I}&\multicolumn{2}{c|}{List-II}\\\hline
(I)&Alpha decay&(P)&${}_{91}\mathrm{Pa}^{230}\to{}_{90}\mathrm{Th}^{230}+\dots$\\\hline
(II)&$\beta^{\ominus}$ decay&(Q)&${}_{30}\mathrm{Zn}^{57}\to{}_{29}\mathrm{Cu}^{56}+\dots$\\\hline
(III)&$\beta^{\oplus}$ decay&(R)&${}_{1}\mathrm{H}^{3}\to{}_{2}\mathrm{He}^{3}+\dots$\\\hline
(IV)&neutron emission&(S)&${}_{53}\mathrm{I}^{137}\to{}_{53}\mathrm{I}^{136}+\dots$\\\hline
&&(T)&${}_{94}\mathrm{Pu}^{239}\to{}_{92}\mathrm{U}^{235}+\dots$\\\hline
\end{tabular}\end{center}
\begin{qoptions}
\opt{A}{I $\to$ T;II $\to$ S;III $\to$ P;IV $\to$ R}
\opt{B}{I $\to$ T;II $\to$ P;III $\to$ Q;IV $\to$ S}
\opt{C}{I $\to$ T;II $\to$ R;III $\to$ P;IV $\to$ S}
\opt{D}{I $\to$ Q;II $\to$ R;III $\to$ S;IV $\to$ P}
\end{qoptions}
\qn{8} List-I contain some optical arrangement and list-II contains deviation produced after all possible optical event in degree then select correct answer using following codes.
\begin{center}\footnotesize\setlength{\tabcolsep}{2pt}
\begin{tabular}{|L{0.7cm}L{4.4cm}|L{0.6cm}L{1.3cm}|}\hline
\multicolumn{2}{|c|}{List-I}&\multicolumn{2}{c|}{List-II}\\\hline
(I)&\qfigi{p1_phy_s2_q08_1}&(P)&0\\\hline
(II)&\qfigi{p1_phy_s2_q08_2}&(Q)&$60^\circ$\\\hline
(III)&\qfigi{p1_phy_s2_q08_3}&(R)&$1.8^\circ$\\\hline
(IV)&\qfigi{p1_phy_s2_q08_4}&(S)&$178.2^\circ$\\\hline
&&(T)&$179.4^\circ$\\\hline
\end{tabular}\end{center}
\begin{qoptions}
\opt{A}{I $\to$ R;II $\to$ P;III $\to$ Q;IV $\to$ R}
\opt{B}{I $\to$ R;II $\to$ Q;III $\to$ Q;IV $\to$ S}
\opt{C}{I $\to$ Q;II $\to$ P;III $\to$ S;IV $\to$ R}
\opt{D}{I $\to$ R;II $\to$ Q;III $\to$ Q;IV $\to$ T}
\end{qoptions}
\qn{9} Here $I$ and $V$ represent current and voltage respectively (steady state for DC and rms for AC) match list-I with list-II using given codes.
\begin{center}\footnotesize\setlength{\tabcolsep}{2pt}
\begin{tabular}{|L{0.7cm}L{4.0cm}|L{0.6cm}L{1.6cm}|}\hline
\multicolumn{2}{|c|}{List-I}&\multicolumn{2}{c|}{List-II}\\\hline
(I)&\qfigi{p1_phy_s2_q09_1}&(P)&$I=0.2$ A\\\hline
(II)&\qfigi{p1_phy_s2_q09_2}&(Q)&$I=0$ A\\\hline
(III)&\qfigi{p1_phy_s2_q09_3}&(R)&$V=20$ Volt\\\hline
(IV)&\qfigi{p1_phy_s2_q09_4}&(S)&$V=0$ Volt\\\hline
&&(T)&$V=10$ Volt\\\hline
\end{tabular}\end{center}
\begin{qoptions}
\opt{A}{I $\to$ P,T;II $\to$ P,R;III $\to$ Q,R;IV $\to$ P,S}
\opt{B}{I $\to$ P,T;II $\to$ P,R;III $\to$ P,S;IV $\to$ P,T}
\opt{C}{I $\to$ P,R;II $\to$ P,S;III $\to$ Q,R;IV $\to$ Q,R}
\opt{D}{I $\to$ P,T;II $\to$ P,R;III $\to$ P,S;IV $\to$ P,S}
\end{qoptions}
\qn{10} List-I contain information about mechanical waves and list-II contain information about velocity of particle of medium corresponding to point $P$ shown. Match list-I with list-II using given codes.
\begin{center}\footnotesize\setlength{\tabcolsep}{2pt}
\begin{tabular}{|L{0.7cm}L{3.5cm}|L{0.6cm}L{2.2cm}|}\hline
\multicolumn{2}{|c|}{List-I}&\multicolumn{2}{c|}{List-II}\\\hline
(I)&\qfigi{p1_phy_s2_q10_1}\par Transverse travelling wave moving in positive $x$-direction with speed 20 m/s&(P)&$v_p=\dfrac{3\pi}{50}$ m/s\\\hline
(II)&\qfigi{p1_phy_s2_q10_2}\par A transverse standing wave oscillating with frequency 50 Hz, having maximum amplitude 2.5 mm&(Q)&$v_p=\dfrac\pi{10}$ m/s\\\hline
(III)&\qfigi{p1_phy_s2_q10_3}\par A longitudinal wave propagating in negative $x$-direction with speed 20 m/s&(R)&moving in positive $x$-direction\\\hline
(IV)&\qfigi{p1_phy_s2_q10_4}\par A longitudinal standing wave oscillating with frequency 50 Hz, having maximum amplitude 2.5 mm.&(S)&moving in positive $y$-direction\\\hline
&&(T)&Insufficient information to find direction of motion\\\hline
\end{tabular}\end{center}
\begin{qoptions}
\opt{A}{I $\to$ Q,S;II $\to$ P,T;III $\to$ Q;IV $\to$ P,T}
\opt{B}{I $\to$ Q,S;II $\to$ P;III $\to$ Q,R;IV $\to$ P,R}
\opt{C}{I $\to$ Q,S;II $\to$ P,T;III $\to$ Q,T;IV $\to$ P}
\opt{D}{I $\to$ P,T;II $\to$ Q,S;III $\to$ Q;IV $\to$ Q,T}
\end{qoptions}
\sect{SECTION-II}
\qn{1} Parallel beam of intensity $I$ incident on a sphere, which absorb 50\% of incident light and reflect remaining light. If radius of sphere is $R$ then radiation force on sphere is given by $\dfrac PQ\,\dfrac{I\left(\pi R^2\right)}c$ then $P+Q$ is (where $P$ and $Q$ are smallest possible natural number) :-
\qn{2} A positron has 0.04 MeV kinetic energy. When it collide with an free electron at rest then they annihilate each other. One of the two photon produced is found to be moving in direction of incident positron. It's energy $x$ MeV. Find $x$. Neglect any relativistic effect. [mass of electron = mass of positron $=\left(\dfrac{0.5}{C^2}\right)$ MeV ]
\qn{3} Radio waves of wavelength 120 m from a galaxy reach a radio telescope by two separate path as shown. As galaxy rises in east over water, the minimum intensity of destructive interference occurs when galaxy is $\theta=37^\circ$ above the horizon. Find the height of the radio telescope disk above the water (in km).
\qfig{p1_phy_s3_q03_1}
\qn{4} Two metal balls of same radius $a$ are located in a homogenous poorly conducting medium with resistivity $\rho$. The resistance of medium between the balls is $\dfrac\rho{Na}$, assuming the separation between them is much greater than radius of the ball find $N$.
\qn{5} In meter bridge setup balanced length is measured by a meter scale having least count of 0.1 cm. When known resistance of $10\ \Omega\pm0.1\ \Omega$ is using then balanced lengths are found to be 40.0 cm. Find \% error in measurement of unknown resistance $x$.
\qfig{p1_phy_s3_q05_1}
\qn{6} \qfig{p1_phy_s3_q06_1}
Taking $R=1$ k$\Omega$ and $E=250$ V find current in wire $ae$ in mA.
\qn{7} A particle with charge $q$ is located at $(-R,0,0)$ and a particle with charge $-2q$ is located at origin then equipotential surface that has zero potential occupy $KR^3$ volume in space. Find $K$.
\qn{8} Two starts are revolving in circular paths about center of mass of system. If there masses are $M$ and $3M$ and separation between them is 4 times separation between the Earth and Sun. If the orbital period is 1 year find ratio of $M$ and mass of Sun.
\partsec{2}{CHEMISTRY}
\sect{SECTION-I(i)}
\qn{1} Which of the following is/are coloured?
\begin{qoptions}
\opt{A}{K$_2$MnO$_4$}
\opt{B}{KO$_2$}
\opt{C}{K$_2$CrO$_4$}
\opt{D}{PbO$_2$}
\end{qoptions}
\qn{2} For phosphinic acid which option is/are correct?
\begin{qoptions}
\opt{A}{The ratio of total number of lone pair electrons \& $\sigma$ bond is 4/5}
\opt{B}{Two atoms are sp$^3$ hybridised}
\opt{C}{It is monobasic acid}
\opt{D}{Value of number of p$\pi$-p$\pi$ bond $\times$ value of number of p$\pi$-d$\pi$ bond = zero}
\end{qoptions}
\qn{3} Select correct statement(s) for the cell with notation :
\par\noindent Ag $|$ AgCl $|$ HCl $||$ AgNO$_3$ $|$ Ag ; $(\Delta S)_{\text{reaction}}=-$ve
\begin{qoptions}
\opt{A}{$E^0_{\text{cell}}>0$}
\opt{B}{$\Delta G^0<0$}
\opt{C}{$E_{\text{cell}}$ decreases by increasing temperature}
\opt{D}{$E_{\text{cell}}$ decreases by diluting anodic solution}
\end{qoptions}
\qn{4} Which option(s) is/are correct for given $P-V$ plot for 2 moles of an ideal gas ? [Given $CA$ is Isothermal process]
\qfig{p1_chem_s1_q04_1}
\begin{qoptions}
\opt{A}{$W_{AB}=-20$ L--atm}
\opt{B}{$\Delta U$ for Cycle $=0$}
\opt{C}{$\Delta S_{\text{Cycle}}=0$}
\opt{D}{$W_{CA}=+13.86$ L atm}
\end{qoptions}
\qn{5} \qfig{p1_chem_s1_q05_1}
Which of the following statement is/are correct for above reaction sequences.
\begin{qoptions}
\opt{A}{Compound (P) can give aldol but do not give tollen's reagent test}
\opt{B}{(P) can give haloform test}
\opt{C}{(R) is a mixture of threo-vicinal dibromide}
\opt{D}{Ozonolysis of (Q) gives two products both of which give positive tollen's reagent test}
\end{qoptions}
\qn{6} Correct option(s) for the following sequence of reactions is(are)
\qfig{p1_chem_s1_q06_1}
\begin{qoptions}
\opt{A}{$Q$ = KNO$_2$, $W$ = LiAlH$_4$}
\opt{B}{$R$ = benzenamine, $V$ = KCN}
\opt{C}{$Q$ = AgNO$_2$, $R$ = phenylmethanamine}
\opt{D}{$W$ = LiAlH$_4$, $V$ = AgCN}
\end{qoptions}
\sect{SECTION-I(ii)}
\qn{7} \begin{center}\footnotesize\setlength{\tabcolsep}{2pt}
\begin{tabular}{|L{0.6cm}L{1.9cm}|L{0.6cm}L{3.4cm}|}\hline
&List-I&&List-II\\\hline
(P)&H$_4$P$_2$O$_7$&(1)&p$\pi$-d$\pi$ bond is present\\\hline
(Q)&P$_4$O$_7$&(2)&P--O--P is present\\\hline
(R)&P$_4$O$_6$&(3)&`P' is sp$^3$ hybridised\\\hline
(S)&P$_4$O$_8$&(4)&Each `P' has 4$\sigma$ \& 1$\pi$ bond\\\hline
&&(5)&Each `P' has only 3$\sigma$ bond\\\hline
\end{tabular}\end{center}
\begin{qoptions}
\opt{A}{P $\to$ 2;Q $\to$ 1;R $\to$ 5;S $\to$ 3}
\opt{B}{P $\to$ 4;Q $\to$ 5;R $\to$ 5;S $\to$ 4}
\opt{C}{P $\to$ 5;Q $\to$ 5;R $\to$ 4;S $\to$ 4}
\opt{D}{P $\to$ 5;Q $\to$ 5;R $\to$ 4;S $\to$ 3}
\end{qoptions}
\qn{8} \begin{center}\footnotesize\setlength{\tabcolsep}{2pt}
\begin{tabular}{|L{0.6cm}L{1.4cm}|L{0.6cm}L{3.3cm}|}\hline
&List-I&&List-II\\\hline
(P)&PbI$_2$&(1)&Water insoluble compound\\\hline
(Q)&AgI&(2)&Yellow colour\\\hline
(R)&Hg$_2$I$_2$&(3)&Green colour\\\hline
(S)&HgI$_2$&(4)&Scarlet red colour\\\hline
&&(5)&White colour\\\hline
\end{tabular}\end{center}
\begin{qoptions}
\opt{A}{P $\to$ 2;Q $\to$ 2;R $\to$ 3;S $\to$ 4}
\opt{B}{P $\to$ 4;Q $\to$ 3;R $\to$ 2;S $\to$ 2}
\opt{C}{P $\to$ 5;Q $\to$ 3;R $\to$ 2;S $\to$ 3}
\opt{D}{P $\to$ 5;Q $\to$ 3;R $\to$ 2;S $\to$ 1}
\end{qoptions}
\qn{9} Answer the following by appropriately matching the lists based on the information given in the list.
\begin{center}\footnotesize\setlength{\tabcolsep}{2pt}
\begin{tabular}{|L{0.6cm}L{3.0cm}|L{0.6cm}L{2.7cm}|}\hline
&List-I&&List-II\\\hline
(P)&Zn$|$\allowbreak Zn$^{2+}$(C)$||$\allowbreak Zn$^{2+}$ (2C)$|$\allowbreak Zn&(1)&Spontaneous cell reaction\\\hline
(Q)&H$_2$(P)$|$\allowbreak HCl(1N)$||$\allowbreak H$_2$SO$_4$(1N)$|$\allowbreak H$_2$(P)&(2)&According to cell representation electric work is produced.\\\hline
(R)&Cu$|$\allowbreak Cu$^{2+}$(0.01M)$||$\allowbreak Ag$^+$(0.1M)$|$\allowbreak Ag&(3)&Concentration of cation in cathodic compartment increases upto equilibrium\\\hline
(S)&Ag$|$\allowbreak AgCl$|$\allowbreak (KCl,0.1M)$||$\allowbreak Ag$^+$(0.01M)$|$\allowbreak Ag&(4)&Concentration cell\\\hline
&&(5)&Concentration of Cl$^-$ in anodic compartment decreases\\\hline
\end{tabular}\end{center}
Correctly match among the following is :

[Given : $E^\circ_{\mathrm{Ag^+/Ag}}=0.8$ V, $E^\circ_{\mathrm{Cu^{+2}/Cu}}=0.34$ V, $K_{sp}$(AgCl) $=10^{-10}$]
\begin{qoptions}
\opt{A}{P $\to$ 1,3;Q $\to$ 3,4;R $\to$ 1,2;S $\to$ 1,2,3,4}
\opt{B}{P $\to$ 1,2,4;Q $\to$ 3,5;R $\to$ 1,3;S $\to$ 1,2,4}
\opt{C}{P $\to$ 1,2,4;Q $\to$ 4;R $\to$ 1,2;S $\to$ 1,2,5}
\opt{D}{P $\to$ 1,2,3;Q $\to$ 3,4;R $\to$ 1,2,5;S $\to$ 1,2,3,5}
\end{qoptions}
\qn{10} \begin{center}\footnotesize\setlength{\tabcolsep}{1pt}
\begin{tabular}{|L{0.6cm}L{3.8cm}|L{0.6cm}L{1.5cm}|}\hline
&List-I&&List-II\\\hline
(P)&\qfigi{p1_chem_s2_q09_1}&(1)&\qfigi{p1_chem_s2_q09_5}\\\hline
(Q)&\qfigi{p1_chem_s2_q09_2}&(2)&\qfigi{p1_chem_s2_q09_6}\\\hline
(R)&\qfigi{p1_chem_s2_q09_3}&(3)&\qfigi{p1_chem_s2_q09_7}\\\hline
(S)&\qfigi{p1_chem_s2_q09_4}&(4)&\qfigi{p1_chem_s2_q09_8}\\\hline
&&(5)&\qfigi{p1_chem_s2_q09_9}\\\hline
\end{tabular}\end{center}
The correct option is :
\begin{qoptions}
\opt{A}{P $\to$ 1;Q $\to$ 2;R $\to$ 3;S $\to$ 4}
\opt{B}{P $\to$ 5;Q $\to$ 4;R $\to$ 3;S $\to$ 1}
\opt{C}{P $\to$ 2;Q $\to$ 5;R $\to$ 3;S $\to$ 1}
\opt{D}{P $\to$ 2;Q $\to$ 4;R $\to$ 5;S $\to$ 3}
\end{qoptions}
\sect{SECTION-II}
\qn{1} For K$_2$Cr$_2$O$_7$,
\par\noindent Total number of $\sigma$-bond $=x$\par\noindent Total number of $\pi$-bond $=y$\par\noindent Total number of lone pair electron $=a$\par\noindent Find the value of $(x+y+a)$
\qn{2} For Borax Crystalline, (Na$_2$B$_4$O$_7\cdot$10H$_2$O)
\par\noindent Total number of $\sigma$ bond $=a$\par\noindent Total number of sp$^3$ hybridised atom $=b$\par\noindent Total number of lone pair on non metal atoms $=c$\par\noindent Total number of sp$^2$ hybridised atom $=d$
\par\noindent\hfill$\dfrac ac\times\dfrac bd=z$\hfill\null\par
Find the value of `z'.
\qn{3} The molar conductivity of a solution of a weak acid HX (0.1 M) is 100 times smaller than the molar conductivity of a solution of a weak acid HY (0.1 M). If $\lambda^0_{X^-}=\lambda^0_{Y^-}$, the difference in their p$K_a$ values, p$K_a$(HX) $-$ p$K_a$(HY), is (consider degree of ionization of both acids to be $\ll1$).
\qn{4} The time for half-life period of a certain reaction. A $\to$ products is 1 h, when the initial concentration of the reactant `A', is 2.0 mol L$^{-1}$. How much time does it take (in hours) for its concentration to come from 0.50 to 0.25 mol L$^{-1}$, if it is a zero order reaction?
\qn{5} 0.1 M aqueous solution of a weak base BOH has pH = 9.5. Find p$K_b$ of the weak base, BOH ?
\qn{6} How many moles of CO$_2$ are released in following reaction :
\qfig{p1_chem_s3_q06_1}
\qn{7} Total number of reactions in which carboxylic acid is obtained as a final product.
\qfig{p1_chem_s3_q07_1}
\qfig{p1_chem_s3_q07_2}
\qn{8} Among the following reaction Identify reactions which give Aromatic product ?
\qfig{p1_chem_s3_q08_1}
\qfig{p1_chem_s3_q08_2}
\qfig{p1_chem_s3_q08_4}
\partsec{3}{MATHEMATICS}
\sect{SECTION-I(i)}
\qn{1} Let $A$, $B$, $C$ be the three sets of complex numbers as defined below $A=\{Z:\operatorname{Im}Z\ge1\}$, $B=\{Z:|Z-2-i|=3\}$ and $C=\left\{Z:\operatorname{Re}\big((1-i)Z\big)=\sqrt2\right\}$. Let `$Z$' be any point in $A\cap B\cap C$ and let `$\omega$' be any point satisfying $|\omega-2-i|<3$, then which of the following may be correct ?
\begin{qoptions}
\opt{A}{$|Z|<3$}
\opt{B}{$|\omega|<4$}
\opt{C}{$|Z|-|\omega|>-6$}
\opt{D}{$|Z|>3$}
\end{qoptions}
\qn{2} Consider the planes $P_1:2x+y+z+4=0$, $P_2:y-z+4=0$ and $P_3:3x+2y+z+8=0$. Let $L_1$, $L_2$, $L_3$ be the lines of intersection of the planes $P_2$ and $P_3$, $P_3$ and $P_1$ \& $P_1$ and $P_2$ respectively. Then,
\begin{qoptions}
\opt{A}{Atleast two of the lines $L_1$, $L_2$ and $L_3$ are non-parallel}
\opt{B}{Atleast two of the lines $L_1$, $L_2$ and $L_3$ are parallel}
\opt{C}{The three planes intersect in a line}
\opt{D}{The three planes form a triangular prism}
\end{qoptions}
\qn{3} A function $y=f(x)$ satisfying the differential equation $\dfrac{dy}{dx}\sin x-y\cos x+\dfrac{\sin^2x}{x^2}=0$ is such that, $y\to0$ as $x\to\infty$, then
\begin{qoptions}
\opt{A}{$\displaystyle\lim_{x\to0}f(x)=1$}
\opt{B}{$\displaystyle\int_0^{\pi/2}f(x)\,dx$ is less than $\dfrac\pi2$}
\opt{C}{$\displaystyle\int_0^{\pi/2}f(x)\,dx$ is greater than unity}
\opt{D}{$f(x)$ is an even function}
\end{qoptions}
\qn{4} Two whole numbers are randomly selected and multiplied. Consider two events $E_1$ and $E_2$ defined as
\par\noindent$E_1$ : Their product is divisible by 5\par\noindent$E_2$ : Unit's place in their product is 5.\par
Which of the following statement(s) is/are correct ?
\begin{qoptions}
\opt{A}{$E_1$ is twice as likely to occur as $E_2$}
\opt{B}{$E_1$ and $E_2$ are disjoint}
\opt{C}{$P(E_2/E_1)=1/4$}
\opt{D}{$P(E_1/E_2)=1$}
\end{qoptions}
\qn{5} Let $f(x)=2x^3-15x^2+36x-23$ \& $g(x)=\begin{cases}\max\{f(t),1\le t\le x\}, & \text{if } 1\le x\le\frac72\\ 12-2x, & \text{if } \frac72<x\le6\end{cases}$. Which of the following is correct about $g(x)$.
\begin{qoptions}
\opt{A}{Continous everywhere}
\opt{B}{Differentiable everywhere}
\opt{C}{has greatest value 5}
\opt{D}{Not differentiable at 2 points}
\end{qoptions}
\qn{6} Let $f:R\to R$ satisfies $f(|f(x)-f(y)|)=f(f(x))-2x^2f(y)+f(y^2)$ where $f(x)$ is not identically zero everywhere. $f(x)$ can not be
\begin{qoptions}
\opt{A}{ONE - ONE}
\opt{B}{ONTO}
\opt{C}{Monotonic}
\opt{D}{MANY-ONE}
\end{qoptions}
\columnbreak
\sect{SECTION-I(ii)}
\qn{7} \begin{center}\footnotesize\setlength{\tabcolsep}{1pt}
\begin{tabular}{|L{0.6cm}L{4.9cm}|L{0.6cm}L{0.7cm}|}\hline
&List-I&&List-II\\\hline
(I)&$A=\begin{bmatrix}2&0&7\\0&1&0\\1&-2&1\end{bmatrix}$ is inverse of $B=\begin{bmatrix}-x&14x&7x\\0&1&0\\1&-4x&-2x\end{bmatrix}$, then the value of $x$ is ?&(P)&3\\\hline
(II)&If $A=\begin{bmatrix}\alpha&2\\2&\alpha\end{bmatrix}$ and $|A^3|=125$ then $\alpha$ is equal to ($\alpha>0$)&(Q)&$\dfrac15$\\\hline
(III)&If $a,b,c$ are three complex numbers such that $a^2+b^2+c^2=0$ and $\Delta=\begin{vmatrix}b^2+c^2&ab&ac\\ab&c^2+a^2&bc\\ac&bc&a^2+b^2\end{vmatrix}=ka^2b^2c^2$&(R)&1\\\hline
(IV)&If $f(\theta)=\begin{vmatrix}\cos^2\theta&\cos\theta\sin\theta&-\sin\theta\\\cos\theta\sin\theta&\sin^2\theta&\cos\theta\\\sin\theta&-\cos\theta&0\end{vmatrix}$ then $f\left(\dfrac{2\pi}3\right)$ is&(S)&4\\\hline
&&(T)&2\\\hline
\end{tabular}\end{center}
\begin{qoptions}
\opt{A}{I $\to$ R;II $\to$ R;III $\to$ T;IV $\to$ T}
\opt{B}{I $\to$ R;II $\to$ P;III $\to$ T;IV $\to$ R}
\opt{C}{I $\to$ Q;II $\to$ P;III $\to$ S;IV $\to$ R}
\opt{D}{I $\to$ Q;II $\to$ P;III $\to$ S;IV $\to$ T}
\end{qoptions}
\qn{8} \begin{center}\footnotesize\setlength{\tabcolsep}{1pt}
\begin{tabular}{|L{0.6cm}L{4.8cm}|L{0.6cm}L{0.7cm}|}\hline
&List-I&&List-II\\\hline
(I)&If the line $\dfrac{x-1}\beta=\dfrac{y-\alpha}{-6}=\dfrac{z-1}{-2}$ entirely lies in the plane $2x+y-6z=0$, then $\alpha+\beta$ is equal to&(P)&$-2$\\\hline
(II)&If planes $2x-y+z=2$, $x-3y+z=\beta$ \& $\alpha x+y+z=5$ have no common point, then $\beta$ cannot be&(Q)&1\\\hline
(III)&If $a,b,c$ are $+$ve real numbers and $ax+by+z=0$, $x+by+cz=0$ \& $ax+y+cz=0$ have infinite common points, then $2abc-ab-bc-ca$ is equal to&(R)&$-1$\\\hline
(IV)&Equation of plane through $(2,1,4)$ and perpendicular to the line $3x-y+z+1=0=5x+y+3z$ is $\alpha x+y+\beta z+\gamma=0$ then $\alpha+\beta+\gamma$ is equal to&(S)&2\\\hline
&&(T)&4\\\hline
\end{tabular}\end{center}
\begin{qoptions}
\opt{A}{I $\to$ Q;II $\to$ R;III $\to$ R;IV $\to$ T}
\opt{B}{I $\to$ R;II $\to$ Q;III $\to$ R;IV $\to$ S}
\opt{C}{I $\to$ Q;II $\to$ R;III $\to$ S;IV $\to$ P}
\opt{D}{I $\to$ R;II $\to$ Q;III $\to$ S;IV $\to$ P}
\end{qoptions}
\qn{9} A showroom has 10 different pairs of expensive shoes. A thief steals 4 shoes randomly from the showroom. Match the following :
\begin{center}\footnotesize\setlength{\tabcolsep}{1pt}
\begin{tabular}{|L{0.6cm}L{3.4cm}|L{0.6cm}L{2.4cm}|}\hline
&List-I&&List-II\\\hline
(I)&The probability that there is no matching pair among the 4 shoes stolen is&(P)&$\dfrac3{19\times17}$\\\hline
(II)&The probability that there is atleast one matching pair among the 4 shoes stolen is&(Q)&$\dfrac{96}{19\times17}$\\\hline
(III)&The probability that there is exactly one matching pair among the 4 shoes stolen is&(R)&$\dfrac{{}^{20}C_4-{}^{10}C_4\cdot2^4}{{}^{20}C_4}$\\\hline
(IV)&The probability that there are two matching pair among the 4 shoes stolen is&(S)&$\dfrac{14\times16}{19\times17}$\\\hline
&&(T)&$\dfrac{16}{19\times17}$\\\hline
\end{tabular}\end{center}
\begin{qoptions}
\opt{A}{I $\to$ R;II $\to$ S;III $\to$ P;IV $\to$ T}
\opt{B}{I $\to$ S;II $\to$ R;III $\to$ T;IV $\to$ Q}
\opt{C}{I $\to$ R;II $\to$ S;III $\to$ Q;IV $\to$ P}
\opt{D}{I $\to$ S;II $\to$ R;III $\to$ Q;IV $\to$ P}
\end{qoptions}
\qn{10} \begin{center}\footnotesize\setlength{\tabcolsep}{1pt}
\begin{tabular}{|L{0.6cm}L{4.6cm}|L{0.6cm}L{1.0cm}|}\hline
&List-I&&List-II\\\hline
(I)&$\displaystyle\int_0^{\pi/4}\ell n\left(\sin^22x\right)dx$&(P)&2\\\hline
(II)&$\displaystyle\int_0^{\pi/2}(x\sin2x-1)\,e^{\cos^2x}\,dx$&(Q)&$-\dfrac\pi2$\\\hline
(III)&If $f\!\left(x+\dfrac\pi{24}\right)+f\!\left(x+\dfrac{5\pi}{24}\right)=f\!\left(x+\dfrac{3\pi}{24}\right)$ and $\displaystyle\int_{\pi/4}^{7\pi/4}f(x)\,dx=3$ then $\displaystyle\int_0^\pi f(x)\,dx$ is equal to&(R)&$-\dfrac\pi2\ell n\,2$\\\hline
(IV)&If $f(x)=\displaystyle\int\dfrac{5x^4+4x^5}{(x^5+x+1)^2}dx$ \& $f(0)=1$ then $f(-1)=$&(S)&$\dfrac\pi2$\\\hline
&&(T)&1\\\hline
\end{tabular}\end{center}
\begin{qoptions}
\opt{A}{I $\to$ Q;II $\to$ Q;III $\to$ T;IV $\to$ T}
\opt{B}{I $\to$ R;II $\to$ Q;III $\to$ P;IV $\to$ P}
\opt{C}{I $\to$ R;II $\to$ S;III $\to$ T;IV $\to$ P}
\opt{D}{I $\to$ Q;II $\to$ S;III $\to$ P;IV $\to$ T}
\end{qoptions}
\sect{SECTION-II}
\qn{1} The value of $\dfrac{4\displaystyle\int_0^\infty x^{10}e^{-x^2}dx}{7\displaystyle\int_0^\infty x^{6}e^{-x^2}dx}$ is
\qn{2} The minimum area bounded by the differentiable function $y=f(x)$ and $y=ax+9$ ($a\in R$) where $f$ satisfies the relation $f(x+y)=f(x)+f(y)+y\sqrt{f(x)}$, $x,y\in R$ and $f'(0)=0$ is
\qn{3} $a,b,c$ are cube roots of $p$($p<0$) then for any permissible real value of $x,y,z$ it is given that $\begin{vmatrix}xa+yb+zc\\xb+yc+za\end{vmatrix}+\left(a_1^2-2b_1^2\right)\omega+([x]+[y]+[z])\omega^2=0$ where $\omega$ is the imaginary cube root of unity and $a_1,b_1$ are prime numbers. Then $[x+a_1]+[y+b_1]+[z]$ is equal to (where $[.]$ denotes G.I.F.)
\qn{4} \par\noindent\hfill$\displaystyle\begin{aligned}\Bigl|\lim_{x\to0}\Bigl[\lim_{n\to\infty}\Bigl\{&\dfrac n{1-n^2x^2}+\dfrac n{4-n^2x^2}\\&+\dots+\dfrac n{n^2-n^2x^2}\Bigr\}\Bigr]\Bigr|\end{aligned}$\hfill\null\par
is equal to
\qn{5} The position vectors of two points $A$ and $C$ are $9\hat i-\hat j+7\hat k$ and $7\hat i-2\hat j+7\hat k$ respectively $\overrightarrow{AB}=4\hat i-\hat j+3\hat k$, $\overrightarrow{CD}=2\hat i-\hat j+2\hat k$, point of intersection of lines $AB$ and $CD$ is $P$. Vector $\overrightarrow{PQ}$ is perpendicular to $\overrightarrow{AB}$ and $\overrightarrow{CD}$ and $PQ=15$. If position vector of $Q$ is $6\hat i-9\hat j-\mu\hat k$ then value of $\mu$ is
\qn{6} If $Z=\dfrac12\left(\sqrt3-i\right)$ then the least positive integral value of `$n$' such that $\left(Z^{101}+i^{109}\right)^{106}=Z^n$ is (where $i^2=-1$)
\qn{7} In a multiple choice question there are four alternative answers of which one or more than one is correct. A candidate will get marks on the question only if he ticks all the correct answer. The candidate ticks the answers at random. The least number of chances he should be given if the probability of getting marks is greater than or equal to $\dfrac13$ is -
\qn{8} Find the distance of the point $\hat i+2\hat j+3\hat k$ from the plane $\vec r\cdot\left(\hat i+\hat j+\hat k\right)=5$ measured parallel to the vector $2\hat i+3\hat j-6\hat k$.
\raggedcolumns
\papersec{2}
\partsec{1}{PHYSICS}
\sect{SECTION-I(i)}
\qn{1} In a hydrogen atom like spectrum if an electron de-excite from higher energy level ($n\gg1$) to a lower energy level then first few frequency of emitted light (Ex : $n$ to $n-1$ or $n$ to $n-2$) are nearly :
\begin{qoptions}
\opt{A}{In Arthmetic progression}
\opt{B}{In Geometric progression}
\opt{C}{In Harmonic progression}
\opt{D}{In ratio $1^2:2^2:3^2$}
\end{qoptions}
\qn{2} A smooth tunnel is made along chord of earth at $\dfrac R2$ from centre of earth ($R$ = radius of earth). A object is released from one end of tunnel. Which of the following is correct graph of variation of contact force $N$ vs distance from centre of tunnel :
\begin{qoptions}
\opt{A}{\qfigi{p2_phy_s1_q02_A}}
\opt{B}{\qfigi{p2_phy_s1_q02_B}}
\opt{C}{\qfigi{p2_phy_s1_q02_C}}
\opt{D}{\qfigi{p2_phy_s1_q02_D}}
\end{qoptions}
\qn{3} A photon of wavelength $\lambda$ is incident on photoelectric plate \& most energetic electron from plate is now incident on target metal of x ray tube. If cut off wavelength of emitted photons from target material is $2\lambda$ then work function of photoelectric plate is :
\begin{qoptions}
\opt{A}{$\dfrac{hc}\lambda$}
\opt{B}{$\dfrac{hc}{2\lambda}$}
\opt{C}{$\dfrac{hc}{4\lambda}$}
\opt{D}{$\dfrac{hc}{8\lambda}$}
\end{qoptions}
\qn{4} If surrounding temperature is slightly fluctuating as shown in figure. A parallel plate capacity with a spacer is design to operate at constant capacitance. If $\alpha_1$ \& $\alpha_2$ are coefficient of linear expansion of plate \& spacer respectively then. (Neglect area of spacer \& it is perfect electric insulator)
\qfig{p2_phy_s1_q04_1}
\begin{qoptions}
\opt{A}{$\alpha_1=\alpha_2$}
\opt{B}{$\alpha_1=2\alpha_2$}
\opt{C}{$2\alpha_1=\alpha_2$}
\opt{D}{$2\alpha_1=3\alpha_2$}
\end{qoptions}
\sect{SECTION-I(ii)}
\qn{5} If there were only one type of charge in the universe then :
\begin{qoptions}
\opt{A}{$\displaystyle\oint_S\vec E\cdot\overrightarrow{dS}\neq0$ on any surface}
\opt{B}{$\displaystyle\oint_S\vec E\cdot\overrightarrow{dS}=0$ if charge is outside the surface}
\opt{C}{$\displaystyle\oint_S\vec E\cdot d\vec S$ could not be defined}
\opt{D}{$\displaystyle\oint_S\vec E\cdot d\vec S=\dfrac q{\varepsilon_0}$ if charges of magnitude $q$ were inside the surface.}
\end{qoptions}
\qn{6} A positive charge $Q$ is uniformly distributed along a non conducting circular ring of radius $R$. A small test charge $q$ of mass $m$ kept on centre of ring. Then :
\qfig{p2_phy_s2_q06_1}
\begin{qoptions}
\opt{A}{If $q>0$ \& is displaced away from center in plane of ring it will be pushed towards center.}
\opt{B}{If $q<0$ \& is displaced away from center in plane of ring then it will never return to center.}
\opt{C}{If $q<0$ \& is displaced along axis then it will perform SHM.}
\opt{D}{If $q<0$ \& is slightly displaced along axis then time period of oscillation is $2\pi\sqrt{\dfrac{mR^3}{kQq}}$.}
\end{qoptions}
\qn{7} A real object is kept perpendicular to principal axis of a lens in air. With respect to different cases of image formation by lens, choose \textbf{CORRECT} option(s) :
\begin{qoptions}
\opt{A}{If magnification is $-2.5$ then image is real \& power of lens is positive.}
\opt{B}{If magnification is 0.5 then image is virtual \& power of lens is negative.}
\opt{C}{If length of image is equal to object then image is real \& power of lens is positive.}
\opt{D}{If length of image is equal to object then image is virtual \& power of lens is negative.}
\end{qoptions}
\qn{8} A cubical region of space is filled with uniform electric \& magnetic fields. A electron \& positron enters with same speed $v$ via opposite faces. At this instant :
\begin{qoptions}
\opt{A}{The electric forces on both particle causes equal acceleration.}
\opt{B}{Acceleration of both particle due to magnetic force are not equal.}
\opt{C}{Rate of change of energy for particles is equal.}
\opt{D}{Motion of COM (center of mass) is determined by only magnetic force.}
\end{qoptions}
\qn{9} In a LCR circuit the average power transfered from a driving source to the circuit is $P=i^2Z\cos\phi$.
\begin{qoptions}
\opt{A}{Here if the power factor $\cos\phi\ge0$ means $P\ge0$.}
\opt{B}{The driving source gives no energy to oscillator in some cases ($P=0$)}
\opt{C}{The driving source cannot removed or syphon out energy out of oscillator ($P<0$)}
\opt{D}{The driving source can take away energy from oscillator.}
\end{qoptions}
\qn{10} Two coherent sources of same intensity in same phase $S_1$ \& $S_2$ of wavelength $\lambda$ are placed as shown. Distance of screen is much larger than $d$.
\qfig{p2_phy_s2_q10_1}
\begin{qoptions}
\opt{A}{If $d=\dfrac{3\lambda}2$ then these will be minimum intensity at $O$.}
\opt{B}{If $d=3\lambda$ then intensity at $O$ is maximum.}
\opt{C}{If $d=\dfrac{11\lambda}6$ then intensity at $O$ is 75\% of maximum intensity.}
\opt{D}{If $d=\lambda$ then minima nearest to $O$ is at $\theta=60^\circ$.}
\end{qoptions}
\sect{SECTION-III}
\qn{1} A galvanometer of coil resistance $1\Omega$ is converted into voltmeter by using $5\Omega$ resistance \& a similar galvanometer is converted into ammeter using $1\Omega$ resistance. Now ammeter \& voltmeter are connected as shown, then ratio of reading of them will be :
\qfig{p2_phy_s3_q01_1}
\qn{2} A ray of light travelling in air is incident at $30^\circ$ on one surface of a slab of variable refractive index. If light travels along path $y=4x^2$ then value of $\mu$ at $y=\dfrac12$ inside the slab is $\dfrac x2$. Find $x$.
\qfig{p2_phy_s3_q02_1}
\qn{3} A loop made of straight edges has 6 corners at $A\,(0,0,0)$ $B\,(L,0,0)$, $C\,(L,L,0)$, $D\,(0,L,0)$, $E\,(0,L,L)$ and $F\,(0,0,L)$ is present in a region of magnetic field $\vec B=\vec B_0\left(\hat i+\hat k\right)$ T. The flux passing through loop is $nB_0L^2$. Find $n$.
\qn{4} Initially at $t=0$ $S_1$ is closed \& $S_2$ is open. After time $RC\ell n2$ $S_1$ is opened \& $S_2$ is closed. At that instant potential difference across capacitor is $V_1$. After time $\dfrac\pi3\sqrt{LC}$ of closing $S_2$ potential difference across $C$ is $V_2$ then $V_1+V_2$ equal to
\qfig{p2_phy_s3_q04_1}
\qn{5} If monochromatic light of $\lambda=5000$\AA\ is incident perpendicular on slits $S_1$ \& $S_2$ ($d=1$mm, $D=1$m). If screen is released from rest under gravity then velocity of $2^{\text{nd}}$ maxima w.r.t $1^{\text{st}}$ maxima on screen at $t=2$ sec (in cm/s). ($O$ is taken as origin.)
\qfig{p2_phy_s3_q05_1}
\qn{6} 3 plates $A$, $B$ \& $C$ each of area 1 m$^2$ are separated by 0.885 mm. Energy stored in system (in $\mu$J) is
\qfig{p2_phy_s3_q06_1}
\qn{7} When a ideal ammeter is put in series in circuit it reads 1 A. When a ideal voltmeter is connected across $X$ it reads 5V. When $A$ \& $B$ are shortened by a conducting wire voltmeter reads 8V across cell. Then ratio of resistance $X$ to internal resistance of battery is
\qfig{p2_phy_s3_q07_1}
\qn{8} There exist a time variable magnetic field centred at origin \& directed towards negative $z$-axis $\left(\dfrac{dB}{dt}=\dfrac{2\text{T}}{\text{s}}\right)$. A conducting rod is kept in region with coordinates of end point as $(1,0,0)$ m \& $(0,2,0)$ m. Then potential difference between end point of rod is :
\partsec{2}{CHEMISTRY}
\sect{SECTION-I(i)}
\qn{1} Which of the following is not paramagnetic.
\begin{qoptions}
\opt{A}{K$_4\!\left[\overset{+2}{\mathrm{Fe}}(\mathrm{CN})_5\mathrm{O}_2\right]$}
\opt{B}{K$_2$MnO$_4$}
\opt{C}{$\overset{+}{\mathrm K}\ \mathrm O_2^-$}
\opt{D}{KMnO$_4$}
\end{qoptions}
\qn{2} In a constant pressure calorimeter 224 ml of 0.1M KOH (aq.) solution is reacted with 50 ml of 0.1M H$_2$SO$_4$ (aq.) solution then increase in temperature of solution will be (assume heat capacity of calorimeter is negligible) -
\par\noindent Given :\par
\noindent Specific heat of solution = 1 cal/gm--K\par
\noindent Density of solution = 1 gm/ml\par
\noindent HCl(aq.) + NaOH(aq.) $\to$ NaCl (aq.) + H$_2$O(l) $\Delta H=-13.7$ kcal/eq.
\begin{qoptions}
\opt{A}{0.5 K}
\opt{B}{1 K}
\opt{C}{2K}
\opt{D}{2.4 K}
\end{qoptions}
\qn{3} For a reaction C$_{\text{(graphite)}}$ + 2H$_{2(g)}\longrightarrow$ CH$_{4(g)}$, $\Delta H^\circ=-75$ KJ at 300K and 1 atm. The standard entropies of C$_{\text{(graphite)}}$, H$_{2(g)}$ and CH$_{4(g)}$ are 6, 130.6 and 186.2 $\dfrac{\text{J}}{\text{K mol}}$ respectively. $\Delta S_{\text{univ}}$ $\left(\text{in}\dfrac{\text{J}}{\text{K}}\right)$ for the reaction will be :-
\begin{qoptions}
\opt{A}{$-81$}
\opt{B}{169}
\opt{C}{$-169$}
\opt{D}{81}
\end{qoptions}
\qn{4} For the given cell at $25^\circ$C, select the correct statement -
\par\noindent Given : $E^0_{\mathrm{Cl_2/Cl^-}}=1.38$ V
\qfig{p2_chem_s1_q04_1}
\begin{qoptions}
\opt{A}{Electrode oxidation potential is negative for anodic compartment.}
\opt{B}{Standard oxidation potential is negative for anodic compartment.}
\opt{C}{Standard Emf of the cell is zero}
\opt{D}{All of these}
\end{qoptions}
\sect{SECTION-I(ii)}
\qn{5} Which of the following will not give chromyl chloride test.
\begin{qoptions}
\opt{A}{CCl$_4$}
\opt{B}{Hg$_2$Cl$_2$}
\opt{C}{C$_2$H$_5$Cl}
\opt{D}{AgCl}
\end{qoptions}
\qn{6} Which of the following will give Tri basic acid on complete hydrolysis?
\begin{qoptions}
\opt{A}{BCl$_3$}
\opt{B}{PCl$_5$}
\opt{C}{AsCl$_3$}
\opt{D}{None}
\end{qoptions}
\qn{7} \qfig{p2_chem_s2_q07_1}
Which option is/are correct for gas (D) and (E).
\begin{qoptions}
\opt{A}{\qfigi{p2_chem_s2_q07_A}}
\opt{B}{Gas `D' has two $\pi$ bonds}
\opt{C}{`E' has two $\pi$ bonds \& shows H-bonding in its liquid}
\opt{D}{Gas `D' is non poisonous}
\end{qoptions}
\qn{8} The correct statement(s) related to colloids is(are)
\begin{qoptions}
\opt{A}{The process of precipitating colloidal sol by adding an electrolyte is called peptization.}
\opt{B}{All colloids may not show Tyndall effect significantly.}
\opt{C}{Surfactants form micelle above critical micelle concentration (CMC), if temperature is above Kraft temperature.}
\opt{D}{Micelles are macromolecular colloids.}
\end{qoptions}
\qn{9} Identify reactions correctly match with their products ?
\begin{qoptions}
\opt{A}{\qfigi{p2_chem_s2_q09_A}}
\opt{B}{\qfigi{p2_chem_s2_q09_B}}
\opt{C}{\qfigi{p2_chem_s2_q09_C}}
\opt{D}{\qfigi{p2_chem_s2_q09_D}}
\end{qoptions}
\qn{10} Reaction in which product can show positive Haloform test :
\begin{qoptions}
\opt{A}{\qfigi{p2_chem_s2_q10_A}}
\opt{B}{\qfigi{p2_chem_s2_q10_B}}
\opt{C}{\qfigi{p2_chem_s2_q10_C}}
\opt{D}{\qfigi{p2_chem_s2_q10_D}}
\end{qoptions}
\sect{SECTION-III}
\qn{1} Find the number of options which are correct for melamine.
\par\noindent(a) Planar (Almost)\par\noindent(b) 15 $\sigma$ bond\par\noindent(c) 3 p$\pi$-p$\pi$ bond are present\par\noindent(d) Resonance takes place\par\noindent(e) Pure sp$^3$ hybridise atoms are = zero\par\noindent(f) All atoms are sp$^2$ hybridised\par\noindent(g) It is used for making unbreakable Utencils
\qn{2} Total number of cation(s) which give white ppt. with K$_4$[Fe(CN)$_6$]
\par\noindent Cu$^{2+}$, Ag$^+$, Ca$^{2+}$, Fe$^{3+}$, Fe$^{2+}$, Zn$^{2+}$, Sr$^{2+}$
\qn{3} The concentration of R in the reaction R $\to$ P was measured as a function of time and the following data is obtained :
\begin{center}\footnotesize\setlength{\tabcolsep}{3pt}
\begin{tabular}{|l|c|c|c|c|}\hline
[R] (molar)&1.0&0.75&0.40&0.10\\\hline
$t$ (min.)&0.0&0.05&0.12&0.18\\\hline
\end{tabular}\end{center}
The order of the reaction is
\qn{4} Liquids \textbf{A} and \textbf{B} form ideal solution over the entire range of composition. At temperature T, equimolar binary solution of liquids \textbf{A} and \textbf{B} has vapour pressure 45 Torr. At the same temperature, a new solution of \textbf{A} and \textbf{B} having mole fractions $x_A$ and $x_B$ respectively, has vapour pressure of 22.5 Torr. If the value of $\dfrac{x_A}{x_B}$ in the new solution is Y. Then the value `Y--10' is :

(Given that the vapour pressure of pure liquid A is 20 Torr at temperature T)
\qn{5} An alcohol (A) on treatment with conc. H$_2$SO$_4$ /$\Delta$ gave an alkene (B). The compound (B), on reacting with Br$_2$/CCl$_4$ and subsequent didehydrobromination with NaNH$_2$(2eq.) produced a compound (C). The compound (C) with dil. H$_2$SO$_4$ in presence of HgSO$_4$ gave a compound `D'. The compound D can also be obtained by oxidation of A by pcc or from dry distillation of calcium acetate. If molecular weight of `C' is Y $\times$ 10, then write the value of Y.
\qn{6} An organic compound A(C$_{10}$H$_{18}$O$_8$) on treatment with excess of CH$_3$COCl gives a fully acetylated product whose molar mass is found to be 518 g/mol. How many hydroxyl functional groups are present in A?
\qn{7} How many total monobrominated products are formed in following reaction sequence :
\par\noindent\hfill$\mathrm{CH_3{-}CH_2{-}CHO}\xrightarrow[\Delta]{\text{dil. NaOH}}(\mathrm A)\xrightarrow[\text{KOH, }\Delta]{\mathrm{NH_2NH_2}}\underset{(\mathrm C_6\mathrm H_{12})}{(\mathrm B)}\xrightarrow[\Delta]{\text{NBS}}\ \text{Mono brominated products}$\hfill\null\par
\qn{8} Scheme 1 and 2 describe the conversion of P to Q and R to S, respectively. Scheme 3 describes the synthesis of T from Q and S. The total number of Br atoms in a molecule of T is \underline{\hspace{1.2cm}}
\par\noindent Scheme 1 :
\qfig{p2_chem_s3_q08_1}
\par\noindent Scheme 2 :
\qfig{p2_chem_s3_q08_2}
\par\noindent Scheme 3 :
\par\noindent\hfill$\mathbf S\xrightarrow[\text{(ii) }\mathbf Q]{\text{(i) NaOH}}\underset{\text{(major)}}{\mathbf T}$\hfill\null\par
\partsec{3}{MATHEMATICS}
\sect{SECTION-I(i)}
\qn{1} Let $f:(0,\infty)\to R$ be a function such that $f(x)=\displaystyle\lim_{n\to\infty}\left(1+\dfrac x{n^2}\right)\left(1+\dfrac{2x}{n^2}\right)\left(1+\dfrac{3x}{n^2}\right)\dots\left(1+\dfrac{nx}{n^2}\right)$. Then correct option is
\begin{qoptions}
\opt{A}{$f'(1)>f(1)$}
\opt{B}{$f(x)<e^x$, $\forall\,x\in(0,\infty)$}
\opt{C}{$f(x)<x^2$, $\forall\,x\in(0,\infty)$}
\opt{D}{$f(c)<\dfrac12$, for some $c$ in $(0,\infty)$}
\end{qoptions}
\qn{2} Let $f:z\to z$ be a function satisfying $f(x)+f(y)=f(x+y)-xy$ and if $f(7)=21$, then which option is correct - (where $z$ is a set of integers)
\begin{qoptions}
\opt{A}{$f(x)$ is an one-one function}
\opt{B}{$f(x)$ is an onto function}
\opt{C}{$f(x)$ is an odd function}
\opt{D}{Inverse function of $f(x)$ does not exist}
\end{qoptions}
\qn{3} $\displaystyle\int\dfrac{x^3-x^2+x-1}{x^5+1}\,dx=f(x)$ and $f(0)=0$, then $f(x)$ is
\begin{qoptions}
\opt{A}{$\dfrac15\ell n\left(\dfrac{1+x^5}{(1+x^3)^5}\right)$}
\opt{B}{$\dfrac15\ell n\left(\dfrac{1+x^5}{1+x^3}\right)$}
\opt{C}{$\dfrac15\ell n\left(\dfrac{1+x^5}{1+x}\right)$}
\opt{D}{$\dfrac15\ell n\left(\dfrac{1+x^5}{(1+x)^5}\right)$}
\end{qoptions}
\qn{4} If $f(x)=\begin{cases}x^n\sin\left(\dfrac1{x^3}\right) & ;\ x\ne0\\ 0 & ;\ x=0\end{cases}$ is continuous but not differentiable at $x=0$, then correct option is
\begin{qoptions}
\opt{A}{$n\in(-1,0)$}
\opt{B}{$n\in(0,1]$}
\opt{C}{$n\in(1,2]$}
\opt{D}{$n\in(2,3]$}
\end{qoptions}
\sect{SECTION-I(ii)}
\qn{5} Let $I(n)=\displaystyle\int_1^n\dfrac{\sin x}{x^2}\,dx$ where $n$ is a natural number, then which of the following options are correct -
\begin{qoptions}
\opt{A}{$I(2)<\ell n2$}
\opt{B}{$I(2)>\ell n2-\dfrac14$}
\opt{C}{$I(2)>\dfrac1{2\sqrt2}$}
\opt{D}{$I(3)<\ell n3$}
\end{qoptions}
\qn{6} Let $P_1$ and $P_2$ be two planes given by
\par\noindent$P_1:10x+15y+12z-60=0$,\par\noindent$P_2:-2x+5y+4z-20=0$.\par
Which of the following straight lines can not be an edge of some tetrahedron whose two faces lie on $P_1$ and $P_2$ ?
\begin{qoptions}
\opt{A}{$\dfrac{x-1}0=\dfrac{y-1}0=\dfrac{z-1}5$}
\opt{B}{$\dfrac{x-6}{-5}=\dfrac y2=\dfrac z3$}
\opt{C}{$\dfrac x{-2}=\dfrac{y-4}5=\dfrac z4$}
\opt{D}{$\dfrac x1=\dfrac{y-4}{-2}=\dfrac z3$}
\end{qoptions}
\qn{7} If $z_1$ is a point on locus given by $|z+8|\le2$ and $z_2$ is a point on locus given by $|z+3|+|z-3|\le10$ then $|z_1-z_2|$ can be
\begin{qoptions}
\opt{A}{0}
\opt{B}{2}
\opt{C}{10}
\opt{D}{18}
\end{qoptions}
\qn{8} $A$ and $B$ are two independent events of an experiment, such that $P(A)=\dfrac13$, $P(B)=\dfrac14$, then
\begin{qoptions}
\opt{A}{$P(A\cup B)=\dfrac12$}
\opt{B}{$P(A\cup\bar B)=\dfrac56$}
\opt{C}{$P(B\cup\bar A)=\dfrac34$}
\opt{D}{$P(A\cap B)=\dfrac1{12}$}
\end{qoptions}
\qn{9} Let $f(x)=\begin{cases}4x^2+1 & x<\frac12\\ 2(\log_ex+1) & \frac12\le x<1\\ x^2+1 & x\ge1\end{cases}$, then which of the following is/are true ?
\begin{qoptions}
\opt{A}{$f(x)$ has two points of extremum}
\opt{B}{$f(x)$ is continuous $\forall\,x\in R$, except at one point}
\opt{C}{$f(x)$ is differentiable $\forall\,x\in R$ except at one point}
\opt{D}{$f(x)$ has exactly one point of minima}
\end{qoptions}
\qn{10} If area bounded by $y^2=8|x|$ and $x^2=3|y|$ is $\Delta$, then $\Delta$ is divisble is
\begin{qoptions}
\opt{A}{4}
\opt{B}{6}
\opt{C}{8}
\opt{D}{12}
\end{qoptions}
\sect{SECTION-III}
\qn{1} If $y=f(x)$ satisfy the differential equation $(1+x^2)f'(x)=x(1-f(x))$, where $f(0)=\dfrac43$, then value of $f\left(\sqrt8\right)-\dfrac19$ is equal to
\qn{2} Let $\beta$ be a real number. Consider the matrix $A=\begin{pmatrix}\beta&0&1\\2&1&-2\\3&1&-2\end{pmatrix}$. If $A^{2022}-(\beta-1)A^{2021}-\beta A^{2020}$ is a singular matrix, then the value of $9\beta$ is \underline{\hspace{1.2cm}}
\qn{3} Let $S$ be the reflection of a point $Q$ with respect to the plane given by $\vec r=-(t+p)\hat i+t\hat j+(1+p)\hat k$ where $t$, $p$ are real parameters and $\hat i$, $\hat j$, $\hat k$ are the unit vectors along the three positive coordinate axes. If the position vectors of $Q$ and $S$ are $10\hat i+15\hat j+20\hat k$ and $\alpha\hat i+\beta\hat j+\gamma\hat k$ respectively, then sum of the digits of $(|\alpha|+|\beta|+|\gamma|)$ is
\qn{4} Let $z$ be a complex number with non-zero imaginary part. If $\dfrac{2+3z+4z^2}{2-3z+4z^2}$ is a real number, then the value of $8|z|^2$ is \underline{\hspace{1.2cm}}.
\qn{5} In a study about a pandemic, data of 900 persons was collected. It was found that
\par\noindent\hspace{1em}190 persons had symptom of fever,\par\noindent\hspace{1em}220 persons had symptom of cough,\par\noindent\hspace{1em}220 persons had symptom of breathing problem,\par\noindent\hspace{1em}330 persons had symptom of fever or cough or both,\par\noindent\hspace{1em}350 persons had symptom of cough or breathing problem or both,\par\noindent\hspace{1em}340 persons had symptom of fever or breathing problem or both,\par\noindent\hspace{1em}30 persons had all three symptoms (fever, cough and breathing problem).\par
If a person is chosen randomly from these 900 persons and the probability that the person has at least two symptoms is $p$, then the value of $10p$ is
\qn{6} In $\Delta ABC$, let $B(3,3,-9)$ and $AB=AC$ and altitude $AD$ is $\dfrac{x-1}1=\dfrac{y-2}2=\dfrac{z+3}3$. If co-ordinates of $C$ are $(\alpha,\beta,\gamma)$, then $|\alpha|+|\beta|+|\gamma|$ is
\qn{7} Let $y=f(x)$ be a curve such that $y$ intercept of tangent at any point is equal to cube of abscissa of point of tangency. If the curve passes through $\left(1,\dfrac12\right)$, then absolute value of $f(2)$ is
\qn{8} The absolute value of $2019\displaystyle\int_0^{\pi/2}(\cos x)^{2018}\cdot\sin(2020x)\,dx$
\end{multicols}
\end{document}
